\documentclass[11pt,a4paper]{article}

\usepackage[margin=2.25cm]{geometry}
\usepackage[T1]{fontenc}
\usepackage[utf8]{inputenc}
\usepackage{lmodern}
\usepackage{microtype}
\usepackage{amsmath,amssymb,amsthm,mathtools,bm}
\usepackage{booktabs,array}
\usepackage{xcolor}
\usepackage{enumitem}
\usepackage{etoolbox}
\usepackage{hyperref}

\AtBeginEnvironment{thebibliography}{\fontsize{8.5}{10}\selectfont\setlength{\itemsep}{0pt}}

\hypersetup{
  colorlinks=true,
  linkcolor=blue!55!black,
  citecolor=blue!55!black,
  urlcolor=blue!55!black,
  pdftitle={Strictly Scalable Exterior Decoders for Quantum Lists: Exact Full-Spark Widths and Fixed-Probe Weyl Bayes Curves},
  pdfauthor={Lluis Eriksson}
}

\newtheorem{theorem}{Theorem}[section]
\newtheorem{proposition}[theorem]{Proposition}
\newtheorem{corollary}[theorem]{Corollary}
\newtheorem{lemma}[theorem]{Lemma}
\newtheorem{remark}[theorem]{Remark}
\newtheorem{example}[theorem]{Example}

\newcommand{\C}{\mathbb{C}}
\newcommand{\I}{\mathbb{I}}
\newcommand{\Tr}{\operatorname{Tr}}
\newcommand{\supp}{\operatorname{supp}}
\newcommand{\cone}{\operatorname{cone}}
\newcommand{\rank}{\operatorname{rank}}
\newcommand{\Span}{\operatorname{span}}
\newcommand{\ket}[1]{\lvert #1\rangle}
\newcommand{\bra}[1]{\langle #1\rvert}
\newcommand{\proj}[1]{\ket{#1}\!\bra{#1}}

\title{\bfseries Strictly Scalable Exterior Decoders for Quantum Lists:\\
Exact Full-Spark Widths and Fixed-Probe Weyl Bayes Curves}
\author{Lluis Eriksson}
\date{August 24, 2026}

\begin{document}
\maketitle

\begin{abstract}
A quantum list measurement succeeds when its classical output contains the
prepared label.  General zero-error feasibility is already characterized by
the learning width, equivalently the factor width of the Gram matrix.  We use
a physical-space form of that known criterion to close an exact
realization-sensitive branch and then go beyond feasibility.  For every
full-spark ensemble of $N$ pure-state rays spanning $\C^r$ that is
\emph{strictly scalable}---its rays admit nonzero tight representatives---the
minimum list size for zero error is
\[
  \ell_{\min}=N-r+1.
\]
The converse is forced by full spark: a smaller-list effect must annihilate
$r$ spanning states and hence vanish.  Attainment is constructive.  Hodge
duals of all $(r-1)$-fold frame wedges form a tight frame, and their rank-one
projectors give a POVM whose output is the complementary list.  The
factor-width criterion becomes an annihilator cone in the physical Hilbert
space; standard conic Carath\'eodory then gives a realization with at most
$r^2$ outcomes.  Below threshold, a smallest-eigenvalue functional gives a
strict Bayes-error floor for every full-spark ensemble and remains positive
under explicit operator-norm perturbations.  Applied to a flat
Schmidt-rank-$r$ probe of the complete $d$-dimensional Weyl-channel
ensemble, the theorem yields the exact fixed-probe frontier
$\ell_{\min}=d-r+1$.  This strictly exceeds both the support-matroid threshold
$\lceil d/r\rceil$ and the global summed-projector obstruction on infinite
families.  An arithmetic-support probe with the same Schmidt rank has a
different exact threshold.  More strongly, when $r\mid d$, a dimension
converse and that arithmetic construction close the optimization over
\emph{all} pure Schmidt-rank-$r$ probes:
$\min_\psi\ell_{\min}(\psi)=d/r$.  For the consecutive rank-two probe we
also determine the complete one-shot Bayes curve,
\[
 P_{\rm succ}^{\star}(d,\ell)
 =\frac{\ell+\sin(\pi\ell/d)/\sin(\pi/d)}{d},
 \qquad 1\le\ell\le d.
\]
The nondivisible probe optimum, adaptive channel testers, and asymptotic
capacity remain outside scope.
\end{abstract}

\section{Introduction}

Quantum state exclusion asks a measurement to certify hypotheses that did not
produce the observed system.  Its list formulation returns a set of candidate
labels and succeeds when the true label lies in that set.  General semidefinite
formulations and optimality conditions are known
\cite{Bandyopadhyay2014}; antidistinguishability has exact algebraic and
Gram-matrix criteria in several settings
\cite{HeinosaariKerppo2018,RussoSikora2023,JohnstonRussoSikora2025}, and
finite group orbits admit exact one-state-exclusion solutions
\cite{DiebraEtAl2026}.  Recent work also studies asymptotic zero-error list
capacity for pure-state classical--quantum channels
\cite{DalaiGirardiLami2026}.  Two-out-of-four state elimination has both an
exact finite analysis and an optical realization
\cite{CrickmoreEtAl2020,WebbEtAl2023}; the multiple-exclusion task itself is
therefore not new here.

Most directly, Johnston, Lovitz, Russo, and Sikora define the same finite
candidate-list task as $k$-learnability and prove that its minimum perfect
list size, the \emph{learning width}, equals the factor width of the state
Gram matrix \cite{JohnstonLovitzRussoSikora2025}.  Their result subsumes the
general feasibility question.  We do not claim a new criterion for arbitrary
ensembles.  Our contribution is an exact evaluation of that width on the
strictly scalable full-spark branch, a weighted exterior-power factorization
that gives the POVM explicitly, and exact
fixed- and optimized-probe Weyl consequences including a complete rank-two
Bayes curve.  We also record a simple perturbatively stable spectral
subthreshold bound.

The present question is deliberately different and finite: given one copy of
one of $N$ prescribed pure states, what is the smallest output list that can
contain the label with certainty?  This question separates two layers.
Linear dependence controls how many states a nonzero effect can annihilate,
while positivity decides whether sufficiently many such effects can resolve
the identity.  A support matroid captures the first layer but not the second
\cite{ErikssonMatroidUnion2026}.  We identify a broad branch where both layers
close exactly.

Let $g_1,\ldots,g_N\in\C^r$ be unit state vectors.  We call their ray
ensemble strictly scalable when positive numbers $s_i$ exist such that
$f_i=s_i g_i$ form a tight frame \cite{KutyniokEtAl2013}.  This is broader than requiring the
physical state vectors themselves to form a unit-norm tight frame, while
remaining invariant under the choice of nonzero ray representatives.
Tightness of the scaled representatives supplies a positive identity
resolution.  Full spark---every $r$ rays are independent---maximizes
resistance to erasure and is a standard frame-theory notion
\cite{AlexeevCahillMixon2012,CasazzaKutyniok2013}.  Neither property
alone determines our threshold.  Together they do:
\[
  \boxed{\ell_{\min}=N-r+1.}
\]
The construction uses exterior algebra.  For every $(r-1)$-subset $E$ of
labels, the Hodge dual $w_E$ of $\bigwedge_{i\in E}f_i$ is orthogonal exactly
to those $r-1$ frame vectors.  Exterior powers of the tight frame operator
show that the projectors $\proj{w_E}$ resolve the identity after one common
normalization.  The measurement reports $[N]\setminus E$.

The cross-product tightness identity itself is classical in real frame
geometry \cite{Ivanov2021}; full-spark harmonic frames are likewise classical
\cite{AlexeevCahillMixon2012}.  We do not claim either ingredient as new.  The
contribution is their exact synthesis into a quantum list decoder, the matching
full-spark converse, a physical-space realization of the known factor-width
criterion, and the resulting Weyl frontiers.

Our main results are:
\begin{enumerate}[leftmargin=2em]
\item The known learning-width/factor-width criterion has an equivalent
physical-space form: the identity lies in a cone generated by projectors onto
vectors with sufficiently many frame zeros (Theorem~\ref{thm:cone}).  This
form also yields the standard $r^2$ conic-Carath\'eodory realization.
\item Every strictly scalable full-spark ensemble of $N$ pure-state rays in
$\C^r$ has exact perfect list threshold $N-r+1$, attained by an explicit
weighted Hodge POVM
(Theorem~\ref{thm:headline}).
\item For arbitrary nonnegative priors and density operators, the simple
spectral quantity $\gamma_\ell$ gives the universal error floor
$P_{\rm err}\ge r\gamma_\ell$.  It is strictly positive below the full-spark
threshold when every prior is positive, and is stable under operator-norm
perturbations
(Theorem~\ref{thm:robustfloor}).
\item For the fixed flat consecutive-support Schmidt-rank-$r$ probe of all
$d^2$ Weyl channels, perfect decoding is possible exactly when
$\ell\ge d-r+1$
(Theorem~\ref{thm:weyl}).
\item For $r=2$, every subthreshold optimum is explicit:
$P_{\rm succ}^{\star}(d,\ell)=
[\ell+\sin(\pi\ell/d)/\sin(\pi/d)]/d$
(Theorem~\ref{thm:ranktwocurve}).
\item If $r\mid d$, optimizing over every pure Schmidt-rank-$r$ probe gives
the exact global value $\min_\psi\ell_{\min}(\psi)=d/r$
(Corollary~\ref{cor:weylopt}).
\end{enumerate}

The hypotheses and resource claims are intentionally narrow.  Full spark and
strict scalability are separately necessary for the stated universal theorem, as
counterexamples show.  The global Weyl optimum is closed only in the divisor
branch $r\mid d$; the nondivisible branch, adaptive uses, and indefinite
causal order are not optimized here.

\section{Learning width, factor width, and annihilator cones}

Let $[N]=\{1,\ldots,N\}$ and let $g_i\in\C^r$ be unit vectors spanning
$\C^r$.  Write $\rho_i=\proj{g_i}$.  A list POVM of size at most $\ell$ is a
family
\[
  \{M_L\succeq0:L\subseteq[N],\ |L|\le\ell\},\qquad
  \sum_L M_L=\I_r.
\]
Zero effects and unused lists are allowed.  This is the weak exclusion
convention; a strong convention requiring every listed subset to occur with a
nonzero effect is a different problem \cite{StrattonHsiehSkrzypczyk2024}.
It is perfect if
\begin{equation}\label{eq:perfect}
  \Tr(\rho_iM_L)=0\qquad\text{whenever }i\notin L.
\end{equation}
For any prior with full support, this is equivalent to unit average success.
Lists of smaller cardinality may be padded, so ``at most $\ell$'' and
``exactly $\ell$'' have the same feasibility threshold.

The \emph{learning width} is the least such feasible $\ell$.  If
$G=[\langle g_i,g_j\rangle]_{i,j=1}^N$ is the Gram matrix, its factor width is
the least $\ell$ for which
\begin{equation}\label{eq:factorwidth}
  G=\sum_a\proj{u_a},
  \qquad |\supp u_a|\le\ell.
\end{equation}
Johnston et al. prove that learning width equals factor width
\cite{JohnstonLovitzRussoSikora2025}.  The next theorem records the equivalent
physical-space form needed for the exterior construction; the equivalence
itself is not claimed new.

For $v\ne0$, define its zero set with respect to the frame
\[
  Z(v)=\{i\in[N]:\langle g_i,v\rangle=0\}
\]
and the admissible annihilator set
\[
  \mathcal V_\ell=\{v\in\C^r\setminus\{0\}:|Z(v)|\ge N-\ell\}.
\]

\begin{theorem}[Factor width in physical-space form]\label{thm:cone}
A perfect list decoder of size at most $\ell$ exists if and only if
\begin{equation}\label{eq:cone}
  \I_r\in\cone\{\proj v:v\in\mathcal V_\ell\}.
\end{equation}
Equivalently, the Gram matrix has factor width at most $\ell$.
If it exists, one can choose a perfect decoder with at most $r^2$ nonzero
outcomes.
\end{theorem}

\begin{proof}
Suppose first that
$\I_r=\sum_{a=1}^m c_a\proj{v_a}$ with $c_a>0$ and
$v_a\in\mathcal V_\ell$.  Set $M_a=c_a\proj{v_a}$ and report
$L_a=[N]\setminus Z(v_a)$.  Then $|L_a|\le\ell$, the effects sum to the
identity, and $i\notin L_a$ implies $\langle g_i,v_a\rangle=0$.  Hence
Eq.~\eqref{eq:perfect} holds.

Conversely, let $\{M_L\}$ be perfect.  Positivity gives
\[
  \langle g_i,M_Lg_i\rangle=0
  \quad\Longrightarrow\quad M_L^{1/2}g_i=0.
\]
Thus every eigenvector with positive eigenvalue in a spectral decomposition of
$M_L$ is orthogonal to all $g_i$ with $i\notin L$, hence belongs to
$\mathcal V_\ell$.  Summing all spectral decompositions proves
Eq.~\eqref{eq:cone}.

Finally, Hermitian $r\times r$ matrices form a real vector space of dimension
$r^2$.  Conic Caratheodory reduces any representation of $\I_r$ in
Eq.~\eqref{eq:cone} to at most $r^2$ generators.  The first part of the proof
turns those generators into a perfect decoder.

For completeness, let $V:\C^N\to\C^r$ have columns $g_i$, so $G=V^*V$.
A cone decomposition $\I_r=\sum_a c_a\proj{v_a}$ maps to
$u_a=\sqrt{c_a}V^*v_a$, giving Eq.~\eqref{eq:factorwidth}; the support of
$u_a$ is the reported list.  Conversely, every factor $u_a$ in
Eq.~\eqref{eq:factorwidth} lies in $\operatorname{ran}G=\operatorname{ran}V^*$,
so write $u_a=V^*v_a$.  Then
$V^*(\sum_a\proj{v_a}-\I_r)V=0$.  Since $V$ is surjective, the middle
operator vanishes.  This is the coordinate translation of the known
factor-width criterion.
\end{proof}

Theorem~\ref{thm:cone} is the realization-sensitive, physical-space form of
the learning-width criterion.  The zero sets are flats of
the represented support matroid, but the inclusion of $\I_r$ in their
projector cone depends on angles and positivity.  For $\ell=N-1$ it is closely
related to known rank-one descriptions of antidistinguishing POVMs
\cite{HeinosaariKerppo2018}; its role here is to connect sparse Gram factors
to physical annihilators and to compress the exterior decoder below.

\section{The exact scalable full-spark threshold}\label{sec:headline}

A nonzero frame $F=(f_i)_{i=1}^N\subset\C^r$ is tight if
\begin{equation}\label{eq:tight}
  \sum_{i=1}^N\proj{f_i}=\alpha\I_r,
  \qquad \alpha=\frac1r\sum_{i=1}^N\|f_i\|^2>0.
\end{equation}
It is full spark if every $r$ frame vectors are linearly independent.
For the associated physical states we normalize only at the end,
\begin{equation}\label{eq:rays}
  g_i=\frac{f_i}{\|f_i\|},\qquad \rho_i=\proj{g_i}.
\end{equation}
Thus the hypothesis concerns the rays: equivalently, the unit vectors $g_i$
form a strictly scalable full-spark ensemble, with positive scaling weights
$s_i=\|f_i\|$.  Individual normalization need not preserve tightness; the
representatives in Eq.~\eqref{eq:tight} are the data used by the exterior
construction.

For $E=\{i_1,\ldots,i_{r-1}\}\subset[N]$, let
\[
  f_E=f_{i_1}\wedge\cdots\wedge f_{i_{r-1}}
  \in\bigwedge^{r-1}\C^r.
\]
Fix an antiunitary Hodge identification
$\star:\bigwedge^{r-1}\C^r\to\C^r$ and put $w_E=\star f_E$.  Equivalently,
$w_E$ is the cofactor vector characterized up to the fixed Hodge convention by
\begin{equation}\label{eq:orth}
  \langle f_i,w_E\rangle=0\quad(i\in E),
  \qquad \|w_E\|^2=\det[\langle f_i,f_j\rangle]_{i,j\in E}.
\end{equation}

\begin{lemma}[Exterior tightness]\label{lem:exterior}
If $F$ obeys Eq.~\eqref{eq:tight}, then
\begin{equation}\label{eq:exterior}
  \sum_{E\in\binom{[N]}{r-1}}\proj{w_E}
  =\alpha^{r-1}\I_r.
\end{equation}
\end{lemma}

\begin{proof}
Let $T:\C^N\to\C^r$ be the synthesis operator $Te_i=f_i$.  Then
$TT^*=\alpha\I_r$.  Apply the $(r-1)$st exterior-power functor:
\[
 (\bigwedge^{r-1}T)(\bigwedge^{r-1}T)^*
 =\bigwedge^{r-1}(TT^*)
 =\alpha^{r-1}\I_{\wedge^{r-1}\C^r}.
\]
The columns of $\bigwedge^{r-1}T$ are the $f_E$.  Conjugation by the
antiunitary Hodge map preserves rank-one sums and sends $f_E$ to $w_E$,
which gives Eq.~\eqref{eq:exterior}.
\end{proof}

The identity is a complex version of the classical cross-product tight-frame
construction \cite{Ivanov2021}.  We use it as a POVM normalization, not as a
new frame-theory claim.

\begin{theorem}[Strictly scalable full-spark list threshold]\label{thm:headline}
Let $F=(f_i)_{i=1}^N$ be a nonzero tight full-spark frame spanning $\C^r$,
with $N>r$, and let $\rho_i$ be the normalized ray states in
Eq.~\eqref{eq:rays}.  The minimum perfect list size is
\begin{equation}\label{eq:threshold}
  \ell_{\min}=N-r+1.
\end{equation}
Equivalently, the Gram matrix of the normalized rays has factor width
$N-r+1$.
At the threshold, an explicit perfect POVM is
\begin{equation}\label{eq:povm}
 M_E=\alpha^{-(r-1)}\proj{w_E},\qquad
 L_E=[N]\setminus E,\qquad |E|=r-1.
\end{equation}
Moreover, a perfect threshold decoder exists using at most $r^2$ of these
complementary lists.
\end{theorem}

\begin{proof}
For the converse, suppose $\ell\le N-r$.  Every allowed list $L$ omits at
least $r$ labels.  Perfectness implies $M_Lg_i=0$, equivalently
$M_Lf_i=0$, for all omitted $i$.  Any
$r$ omitted frame vectors span $\C^r$ by full spark, so $M_L=0$.  This would
hold for every outcome, contradicting $\sum_LM_L=\I_r$.  Therefore
$\ell\ge N-r+1$.

For attainment, Lemma~\ref{lem:exterior} shows that the effects in
Eq.~\eqref{eq:povm} sum to $\I_r$.  Equation~\eqref{eq:orth} and
$g_i\parallel f_i$ show that $M_E$ annihilates every physical state with label
in $E$, exactly the labels omitted from the
reported list.  Hence the decoder is perfect with list size $N-r+1$.

Full spark makes each $(r-1)$-fold wedge nonzero and prohibits a vector from
having $r$ frame zeros.  Consequently every generator of the threshold cone
in Theorem~\ref{thm:cone} is proportional to some $w_E$.  Conic Caratheodory
therefore selects at most $r^2$ of the displayed effects, with possibly
different positive weights, while retaining the identity resolution.
\end{proof}

\begin{proposition}[Constructive $r^2$ compression]\label{prop:compression}
The symmetric Hodge POVM in Eq.~\eqref{eq:povm} can be compressed to at most
$r^2$ nonzero outcomes by finitely many real-nullspace computations and
nonnegative weight updates.  The retained outcomes report the same lists and
remain a perfect decoder.
\end{proposition}

\begin{proof}
Write the current identity resolution as
$\I_r=\sum_{j=1}^m c_jP_j$, where $c_j>0$ and every $P_j$ is one of the Hodge
projectors.  Regard Hermitian matrices as vectors in the real space
$\operatorname{Herm}(r)$ of dimension $r^2$.  If $m>r^2$, compute a nonzero
real vector $\beta$ satisfying $\sum_j\beta_jP_j=0$.  Its coefficients have
both signs: a nonzero positive combination of nonzero positive semidefinite
matrices cannot vanish.  After changing the sign of $\beta$ if necessary,
set
\[
 t=\min_{\beta_j>0}\frac{c_j}{\beta_j},
 \qquad c'_j=c_j-t\beta_j.
\]
Then every $c'_j\ge0$, at least one becomes zero, and
$\sum_jc'_jP_j=\I_r$.  Delete zero coefficients and repeat.  The procedure
terminates with at most $r^2$ effects.  Only weights change, so every retained
effect keeps its original annihilated set and reported list.
\end{proof}

\begin{remark}[What each hypothesis does]
Full spark is used only in the converse and to identify the maximal zero sets.
Tightness is used only to place the identity in the annihilator cone through
Eq.~\eqref{eq:exterior}.  Section~\ref{sec:counterexamples} shows that neither
hypothesis can be deleted from the universal statement.
\end{remark}

\begin{remark}[Weighted physical form]\label{rem:weighted}
If unit state representatives $g_i$ obey
$\sum_i\lambda_i\proj{g_i}=\I_r$ with every $\lambda_i>0$, take
$f_i=\sqrt{\lambda_i}g_i$ and $\alpha=1$.  Then the threshold effects are
the Hodge projectors of the weighted wedges
$\bigwedge_{i\in E}\sqrt{\lambda_i}g_i$.  Hence the theorem is genuinely
projective: unequal representative norms are not physical probabilities, but
a strictly positive scalability certificate used to synthesize the POVM.
The unit-norm tight theorem is the special case $\lambda_i=r/N$.
\end{remark}

The number $\binom{N}{r-1}$ of effects in Eq.~\eqref{eq:povm} is a symmetric
description, not a hardware lower bound.  Proposition~\ref{prop:compression}
turns the standard conic-Caratheodory bound into an explicit finite
compression algorithm.  It requires nullspaces of real matrices with $r^2$
rows; it is not asserted to minimize the number of outcomes or to be
numerically well conditioned without certified linear algebra.

\section{A spectral subthreshold certificate}\label{sec:robust}

The zero-error threshold is discontinuous as a yes/no statement, but its
failure admits a quantitative and perturbatively stable certificate.  Let
$p_i\ge0$, $\sum_i p_i=1$, and let $\rho_i$ be arbitrary density operators on
$\C^r$.  For every allowed list $L$ define the omitted-prior operator
\[
 Q_L=\sum_{i\notin L}p_i\rho_i,
 \qquad
 \gamma_\ell=\min_{|L|\le\ell}\lambda_{\min}(Q_L).
\]

\begin{theorem}[Spectral list-error bound]\label{thm:robustfloor}
For every list POVM of size at most $\ell$,
\begin{equation}\label{eq:robustfloor}
 P_{\rm err}\ge r\gamma_\ell.
\end{equation}
If the states are pure, have full-support priors, and their rays are full
spark, then $\gamma_\ell>0$ whenever $\ell\le N-r$.
Moreover, if $\|\widetilde\rho_i-\rho_i\|_{\rm op}\le\varepsilon_i$, then
\begin{equation}\label{eq:perturbedfloor}
 \widetilde P_{\rm err}
 \ge r\max\!\left\{0,
 \min_{|L|\le\ell}\left[
 \lambda_{\min}(Q_L)-\sum_{i\notin L}p_i\varepsilon_i
 \right]\right\}
 \ge r\max\!\left\{0,
 \gamma_\ell-\sum_i p_i\varepsilon_i\right\}.
\end{equation}
\end{theorem}

\begin{proof}
For a POVM $\{M_L\}$, the probability of omitting the true label is
\[
 P_{\rm err}=\sum_L\Tr(M_LQ_L).
\]
Since $Q_L\succeq\gamma_\ell\I_r$ for every allowed list,
\[
 P_{\rm err}\ge\gamma_\ell\sum_L\Tr M_L
 =r\gamma_\ell.
\]
If $\ell\le N-r$, every $L$ omits at least $r$ labels.  For a full-spark pure
ensemble those omitted rays span $\C^r$; because their priors are positive,
$Q_L$ is positive definite.  There are finitely many lists, so their smallest
eigenvalues have a strictly positive minimum.

Finally,
$\|\widetilde Q_L-Q_L\|_{\rm op}
\le\sum_{i\notin L}p_i\varepsilon_i$.
Weyl's eigenvalue perturbation inequality gives the first lower bound in
Eq.~\eqref{eq:perturbedfloor}; replacing every omitted-list sum by
$\sum_i p_i\varepsilon_i$ gives the second.
\end{proof}

This is the scalar smallest-eigenvalue witness obtained from the standard
quantum decision SDP \cite{Holevo1973,YuenKennedyLax1975}; no novelty is
claimed for that dual mechanism.  Its usefulness here is that it depends on
the actual angles and priors, not only on the support
matroid.  It therefore supplies quantitative information that the zero-error
width discards.  It is a certified lower bound, not generally the exact Bayes
error; Section~\ref{sec:weylcurve} gives a nontrivial family whose complete
optimum can nevertheless be solved.

\section{What support-only obstructions miss}

For a full-spark $N$-frame in $\C^r$, the rank of a subfamily is
$\min\{|A|,r\}$.  The $\ell$-fold support-matroid obstruction therefore
vanishes once
\begin{equation}\label{eq:matroidthreshold}
  \ell\ge\left\lceil\frac Nr\right\rceil.
\end{equation}
This is a necessary support threshold, not a measurement construction
\cite{ErikssonMatroidUnion2026}.  Theorem~\ref{thm:headline} gives the exact
physical threshold $N-r+1$, which may be much larger.

For the unit-norm tight subfamily the common global projector test also becomes
\[
  \sum_i\rho_i=\frac Nr\I_r\preceq\ell\I_r
\]
at Eq.~\eqref{eq:matroidthreshold}.  Thus both coarse tests can be silent while
perfect decoding remains impossible.  We stress ``global'': stronger
subset-wise projector criteria exist in exclusion theory
\cite{StrattonHsiehSkrzypczyk2024}, and no equivalence with them is claimed.

\begin{table}[t]
\centering
\caption{Support threshold versus the exact perfect-list threshold for a
representative unit-norm tight subfamily.}
\label{tab:gaps}
\begin{tabular}{cccc}
\toprule
$N$ & $r$ & support/global threshold $\lceil N/r\rceil$
& exact $N-r+1$\\
\midrule
4 & 2 & 2 & 3\\
6 & 2 & 3 & 5\\
6 & 3 & 2 & 4\\
8 & 4 & 2 & 5\\
10& 5 & 2 & 6\\
\bottomrule
\end{tabular}
\end{table}

For example, keeping $r=2$ and increasing $N$ produces an infinite strict
family: the support obstruction disappears at $\lceil N/2\rceil$, while the
physical threshold is $N-1$.

\section{A fixed-probe Weyl-channel frontier}\label{sec:weyl}

We now turn the frame theorem into an exact communication statement.  Let
$X\ket{x}=\ket{x+1\bmod d}$ and $Z\ket{x}=\omega^x\ket{x}$, where
$\omega=e^{2\pi i/d}$.  One of the $d^2$ equiprobable channels
\[
  \mathcal U_{ab}(\cdot)=X^aZ^b(\cdot)Z^{-b}X^{-a},
  \qquad (a,b)\in\mathbb Z_d^2,
\]
is used once.  Dense coding is the rank-$d$ endpoint
\cite{BennettWiesner1992}.  Partially entangled deterministic dense coding and
its probe dependence are classical
\cite{MozesOppenheimReznik2005,WuCohenSunGriffiths2006}.  We fix the specific
flat consecutive-support Schmidt-rank-$r$ probe
\begin{equation}\label{eq:probe}
  \ket{\Phi_r}=\frac1{\sqrt r}\sum_{x=0}^{r-1}\ket{x}_A\ket{x}_R,
  \qquad 1\le r\le d.
\end{equation}
No optimization over probes is implicit.
The companion record \cite{ErikssonMatroidUnion2026} gives a general
unitary-error-basis list cap and reports the point $(d,r,\ell)=(4,2,2)$ for a
consecutive-support probe.  The theorem below strictly generalizes that
fixed-probe comparison across all $d$ and $r$; it does not re-claim the
companion fixture.

\begin{theorem}[Fixed-probe Weyl frontier]\label{thm:weyl}
For the complete $d$-dimensional Weyl ensemble and the probe
Eq.~\eqref{eq:probe}, perfect one-use list decoding is possible if and only if
\begin{equation}\label{eq:weylthreshold}
  \ell\ge d-r+1.
\end{equation}
At equality, a decoder first reads the shift sector and then applies the Hodge
POVM of Theorem~\ref{thm:headline} to the phase frame in that sector.
\end{theorem}

\begin{proof}
The output vector is
\begin{equation}\label{eq:weyloutput}
  \ket{\psi_{ab}}=\frac1{\sqrt r}
  \sum_{x=0}^{r-1}\omega^{bx}\ket{x+a}_A\ket{x}_R.
\end{equation}
Different $a$ give mutually orthogonal subspaces.  For fixed $a$, the $d$
vectors are unitarily equivalent to
\[
  \ket{\varphi_b}=\frac1{\sqrt r}
  (1,\omega^b,\ldots,\omega^{(r-1)b})^T\in\C^r.
\]
They form a unit-norm tight frame with bound $d/r$.  Every $r\times r$ minor
of the synthesis matrix is a Vandermonde determinant on distinct $d$th roots
of unity, so the frame is full spark
\cite{AlexeevCahillMixon2012}.  Theorem~\ref{thm:headline} supplies a perfect
within-sector list of size $d-r+1$, and the shift measurement identifies $a$
without error.  This proves sufficiency.

For necessity, let $M_L$ be any effect of a perfect global list measurement.
Replacing every effect by its pinching into the orthogonal shift sectors
preserves positivity, completeness, and all probabilities of the states in
Eq.~\eqref{eq:weyloutput}; hence we may take every effect block diagonal.  In
each sector, $L$ contains
at most $\ell$ phase labels.  If $\ell\le d-r$, the effect's block must
annihilate at least $r$ phase states; full spark forces that block to vanish.
Every block vanishes, hence $M_L=0$.  Since all effects would vanish, no POVM
exists.  Thus Eq.~\eqref{eq:weylthreshold} is necessary.
\end{proof}

\subsection{The complete rank-two Bayes curve}\label{sec:weylcurve}

For $r=2$, the phase ensemble in each shift sector is the regular equatorial
$d$-gon
\[
 \ket{\varphi_b}=\frac{\ket0+\omega^b\ket1}{\sqrt2}.
\]

\begin{lemma}[Largest sum of polygon vertices]\label{lem:polygon}
For $0\le k\le d$,
\[
 \max_{S\subseteq\mathbb Z_d,\ |S|=k}
 \left|\sum_{b\in S}\omega^b\right|
 =\frac{\sin(\pi k/d)}{\sin(\pi/d)},
\]
with the right-hand side interpreted as zero for $k=0,d$.  A consecutive
arc attains the maximum.  Endpoint projection ties may give more than one
maximizing arc but do not change the value.
\end{lemma}

\begin{proof}
Let $S$ maximize the modulus and, when its sum is nonzero, let $u$ be its
unit complex direction.  Then
\[
 \left|\sum_{b\in S}\omega^b\right|
 =\sum_{b\in S}\Re(\overline u\omega^b).
\]
Replacing $S$ by the $k$ vertices with largest projection on $u$ cannot
decrease this expression, while the modulus of their sum is at least that
projection.  The $k$ largest projections on any direction form a consecutive
arc of the regular polygon, with only the stated endpoint ambiguity.  A
geometric-series evaluation of such an arc gives the displayed sine ratio
for $1\le k<d$.  When $k>d/2$, the same value also follows by complementing:
the sum over the complement is the negative of the original sum because all
$d$ roots sum to zero.  The cases $k=0,d$ are immediate.
\end{proof}

Minimum-error discrimination of geometrically uniform states and the
associated square-root measurement are classical
\cite{EldarForney2001}.  The following result solves the different
list-valued objective for every list size.

\begin{theorem}[Exact rank-two Weyl list curve]\label{thm:ranktwocurve}
For the complete uniform $d^2$ Weyl ensemble and the consecutive-support
probe $\ket{\Phi_2}$, the optimal one-use success probability with lists of
size at most $\ell$ is
\begin{equation}\label{eq:ranktwocurve}
 P_{\rm succ}^{\star}(d,\ell)
 =\frac{1}{d}\left(
 \ell+\frac{\sin(\pi\ell/d)}{\sin(\pi/d)}\right),
 \qquad 1\le\ell\le d.
\end{equation}
In particular $P_{\rm succ}^{\star}<1$ for $\ell\le d-2$ and equals one for
$\ell=d-1,d$, recovering Theorem~\ref{thm:weyl} at $r=2$.
\end{theorem}

\begin{proof}
Pinch any global measurement into the orthogonal shift sectors as in the
proof of Theorem~\ref{thm:weyl}.  For a set $L\subseteq\mathbb Z_d$ of phase
labels, $|L|\le\ell$, put
\[
 A_L=\sum_{b\in L}\proj{\varphi_b}
 =\frac12\begin{pmatrix}
 |L|&\sum_{b\in L}\omega^{-b}\\
 \sum_{b\in L}\omega^b&|L|
 \end{pmatrix}.
\]
Hence
\[
 \lambda_{\max}(A_L)
 =\frac12\left(|L|+\left|\sum_{b\in L}\omega^b\right|\right).
\]
Lemma~\ref{lem:polygon} shows that the largest possible root sum has magnitude
\begin{equation}\label{eq:rootsum}
 R_{d,\ell}=\frac{\sin(\pi\ell/d)}{\sin(\pi/d)}.
\end{equation}

For the upper bound, write
$L_a=\{b:(a,b)\in L\}$.  If $|L_a|<\ell$, padding it to $\ell$ phase labels
can only increase the positive operator $A_{L_a}$.  Hence each pinched effect
block $M_L^{(a)}$ obeys
\[
 \Tr\!\left(M_L^{(a)} A_{L_a}\right)
 \le \frac{\ell+R_{d,\ell}}{2}\Tr M_L^{(a)}.
\]
Summing over lists and the $d$ two-dimensional shift sectors gives
\[
 P_{\rm succ}
 \le \frac{1}{d^2}\frac{\ell+R_{d,\ell}}{2}
       \sum_{a,L}\Tr M_L^{(a)}
 =\frac{1}{d^2}\frac{\ell+R_{d,\ell}}{2}(2d),
\]
which is Eq.~\eqref{eq:ranktwocurve}.

If $\ell=d$, the single effect $\I_2$ labelled by the full list attains one,
so suppose $\ell<d$.  Choose a consecutive base list $L_0$ and a unit eigenvector
$v_0$ of $A_{L_0}$ at its largest eigenvalue.  Translate both by the cyclic
phase action: $L_t=L_0+t$ and
$v_t=\operatorname{diag}(1,\omega^t)v_0$.  Every $v_t$ is equatorial and
$\sum_{t=0}^{d-1}\proj{v_t}=(d/2)\I_2$, so
$M_t=(2/d)\proj{v_t}$ is a POVM.  It attains the largest eigenvalue for every
translated list.  Applying this covariant list POVM after the projective
shift measurement proves equality in every sector and hence globally.
\end{proof}

The endpoints have a direct meaning:
\[
 r=1:\ \ell_{\min}=d,
 \qquad
 r=d:\ \ell_{\min}=1.
\]
Each additional flat Schmidt coefficient removes exactly one label from the
zero-error list.  Rearranging Eq.~\eqref{eq:weylthreshold} gives the exact
fixed-family memory witness
\[
 r\ge d-\ell+1.
\]
Here $r$ is the Schmidt rank of the specified probe, not a claim about a
general quantum-memory cost.

\begin{remark}[The support pattern matters]\label{rem:arithmeticprobe}
The frontier is not determined by Schmidt rank alone.  If $r\mid d$ and
$q=d/r$, the equally entangled arithmetic-support probe
\[
 \ket{\widetilde\Phi_r}=\frac1{\sqrt r}
 \sum_{t=0}^{r-1}\ket{qt}_A\ket{t}_R
\]
produces mutually orthogonal sectors labelled by $a$ and $b\bmod r$; each
output vector is shared by exactly the $q$ labels with the same residue class.
Measuring those sectors therefore gives the exact threshold
$\ell_{\min}=q=d/r$.  Necessity follows because the $q$ corresponding labels
produce identical output states.  Thus, for example, $(d,r)=(4,2)$ gives
$\ell_{\min}=2$ for this probe but $3$ for Eq.~\eqref{eq:probe}.  Theorem
\ref{thm:weyl} is consequently a support-pattern theorem.  The next
corollary shows that the arithmetic pattern is globally optimal in the
divisor branch.
\end{remark}

\begin{corollary}[Optimized Weyl divisor branch]\label{cor:weylopt}
Suppose $r\mid d$.  Among all pure probes of Schmidt rank exactly $r$ for the
complete uniform $d^2$ Weyl ensemble,
\begin{equation}\label{eq:globalweyl}
  \min_{\psi:\,\operatorname{SR}(\psi)=r}\ell_{\min}(\psi)=\frac dr.
\end{equation}
The arithmetic-support probe in Remark~\ref{rem:arithmeticprobe} is optimal.
The lower bound specializes the support-dimension mechanism already used for
complete unitary error bases in \cite{ErikssonMatroidUnion2026}; the new point
here is its matching arithmetic-support construction for every divisor pair.
\end{corollary}

\begin{proof}
Fix any such probe and let $\mathcal K$ be the span of its $d^2$ pure output
vectors.  Since the channel output has dimension $d$ and the probe's reference
support has dimension $r$, $D=\dim\mathcal K\le dr$.  Compress any perfect
list POVM to $\mathcal K$, so that $\sum_LM_L=\I_{\mathcal K}$.  Writing
$\rho_{ab}$ for the output states and using
$\sum_{(a,b)\in L}\rho_{ab}\preceq |L|\I_{\mathcal K}$ gives
\begin{align*}
  1
  &=\frac1{d^2}\sum_L
    \Tr\!\left(M_L\sum_{(a,b)\in L}\rho_{ab}\right)\\
  &\le \frac{\ell}{d^2}\sum_L\Tr M_L
   =\frac{\ell D}{d^2}
  \le\frac{\ell r}{d}.
\end{align*}
Thus every rank-$r$ probe needs $\ell\ge d/r$.  The arithmetic-support probe
attains equality by Remark~\ref{rem:arithmeticprobe}.
\end{proof}

\section{Necessity of strict scalability and full spark}\label{sec:counterexamples}

\subsection{Tight but not full spark}

In $\C^2$, take
\[
 f_1=f_2=e_1,\qquad f_3=f_4=e_2.
\]
This is a unit-norm tight frame but not full spark.  The basis measurement
reports $\{1,2\}$ or $\{3,4\}$, so $\ell=2$ is perfect, below
$N-r+1=3$.  Full spark cannot be removed from the converse.

\subsection{Full spark but not strictly scalable}

Take the three qubit vectors
\[
 f_1=\ket0,\qquad
 f_2=\frac{3\ket0+\ket1}{\sqrt{10}},\qquad
 f_3=\frac{3\ket0-\ket1}{\sqrt{10}}.
\]
Every pair is independent, so the frame is full spark.  All three Bloch
vectors have strictly positive $z$ component.  A list-two effect must omit at
least one label and is therefore supported on the line orthogonal to the
omitted state.  Its Bloch vector has strictly negative $z$ component.  No
positive combination of such nonzero effects can sum to the identity, whose
Bloch vector is zero.  Hence list two is impossible although
$N-r+1=2$.  The same hemisphere argument shows that no strictly positive
rescaling of these rays is tight: every weighted Bloch sum still has positive
$z$ component.  Strict scalability cannot be removed from attainment.  This is the qubit
cone criterion of \cite{HeinosaariKerppo2018} in the present list language.

\subsection{Other boundaries}

Zero-prior labels must be removed before defining $N$.  A nonzero effect
cannot annihilate a full-rank state, so replacing frame vectors by arbitrary
mixed states is invalid.  Near-tight frames do not inherit exact zero error by
continuity: feasibility is a closed conic condition with exact orthogonality
constraints.  Theorem~\ref{thm:robustfloor} instead supplies a quantitative
floor under bounded state perturbations, while
Theorem~\ref{thm:ranktwocurve} gives an exact subthreshold curve only for the
specified rank-two Weyl family.

\section{Relation to prior work and scope}

Table~\ref{tab:prior} states the contribution boundary.  General exclusion
SDPs, antidistinguishability, group-covariant one-state exclusion, exterior
tightness, full-spark harmonic frames, and the equality between learning
width and Gram-matrix factor width, as well as the notion of factor-width
rank, all predate this work.
In particular, factor-width rank records the minimum number of sparse
rank-one terms \cite{JohnstonMoeinPlosker2025}; the $r^2$ bound below is the
standard conic-Carath\'eodory consequence in the real vector space of
Hermitian operators on the $r$-dimensional support.  In
particular, the group-generated solution of \cite{DiebraEtAl2026} concerns an
outcome that excludes its matching single label.  Our threshold decoder
instead associates an outcome with $r-1$ simultaneously excluded labels and
solves the smallest retained list over an arbitrary strictly scalable
full-spark ray ensemble.
The channel-capacity results of \cite{DalaiGirardiLami2026} are asymptotic in
block length; our result is a one-shot, fixed-codebook feasibility law.
The companion record \cite{ErikssonMatroidUnion2026} already proves a general
unitary-error-basis list cap and the single $(d,r,\ell)=(4,2,2)$
subthreshold fixture.  Our formulas strictly generalize that reported point;
we do not claim it again as an independent contribution.  The new Weyl
content here is the exact all-$(d,r)$ zero-error threshold for the specified
consecutive-support probe, its Hodge decoder, and the exact optimization over
all rank-$r$ pure probes when $r\mid d$, together with the complete
rank-two list-valued Bayes curve.

\begin{table}[t]
\centering
\small
\caption{Contribution boundary.  ``Used'' means explicitly credited rather
than claimed as new.}
\label{tab:prior}
\begin{tabular}{@{}>{\raggedright\arraybackslash}p{0.29\linewidth}
                    >{\raggedright\arraybackslash}p{0.30\linewidth}
                    >{\raggedright\arraybackslash}p{0.32\linewidth}@{}}
\toprule
Topic & Status here & New synthesis claimed here\\
\midrule
Learning width $=$ Gram factor width & Used
& Exact evaluation $N-r+1$ on the scalable full-spark branch\\
Exclusion SDP and optimality & Used
& Exact scalable full-spark list threshold\\
Pure-state annihilator/factor cones and factor-width rank & Used; the $r^2$
compression is a standard Carath\'eodory corollary
& Weighted Hodge factorization and exact full-spark width\\
Cross-products/exterior powers of tight frames & Used
& Hodge projectors as a threshold POVM\\
Full-spark Fourier/Vandermonde frames & Used
& Exact consecutive-support Weyl frontier\\
Dense coding and Weyl channels & Used
& Fixed-probe frontier, global divisor optimum, and rank-two Bayes curve\\
Geometrically uniform state discrimination & Used
& Exact list-valued regular-polygon curve\\
\bottomrule
\end{tabular}
\end{table}

We make no ``first ever'' priority claim.  The general feasibility criterion
is explicitly credited to \cite{JohnstonLovitzRussoSikora2025}.  Targeted
comparison with the primary sources above did not locate the scalable
all-frame evaluation $N-r+1$, the complementary weighted Hodge decoder, the
consecutive-support Weyl frontier $d-r+1$, the global divisor-branch optimum
$d/r$, or the exact rank-two list curve.
The novelty claim is restricted to those evaluations and constructions, not
to factor width, exclusion, exterior algebra, covariant discrimination, or
dense coding individually.

\section{Reproducibility}

The accompanying lightweight NumPy replay checks seven finite layers:
\begin{enumerate}[leftmargin=2em]
\item harmonic frames are tight and full spark for a grid of small $(N,r)$;
\item exterior Hodge effects reproduce $\alpha^{r-1}\I$, including a
non-unit-norm tight full-spark fixture;
\item nullspace elimination compresses a finite Hodge POVM to at most $r^2$
effects while retaining the identity resolution;
\item consecutive- and arithmetic-support Weyl fixtures agree with the closed
formulas, including the divisor-branch optimum;
\item the rank-two Weyl Bayes formula agrees with explicit covariant POVMs for
a grid of dimensions and list sizes;
\item the spectral subthreshold floor and its perturbation inequality hold on
finite harmonic-frame fixtures;
\item both hypothesis counterexamples behave as stated.
\end{enumerate}
The replay is diagnostic evidence, not a proof of the all-dimensional
theorems.  Those follow from the exterior-power, full-spark, and Weyl
arguments in Sections~\ref{sec:headline} and~\ref{sec:weyl}.  No
formal-proof-assistant certification, peer review, or independent scientific
validation is claimed.

\section{Conclusion}

For quantum list decoding, linear independence alone is an obstruction and
positivity is the construction problem.  Strictly scalable full-spark ray
ensembles make the two meet exactly: full spark forces
$\ell\ge N-r+1$, while the
$(r-1)$st exterior power of tightness builds the matching POVM.  The known
factor-width criterion, in physical-space annihilator form, explains why this
construction is physical; standard conic Carath\'eodory explains the
$r^2$-outcome compression.

The Weyl application turns the abstract threshold into a sharp resource
frontier for a fixed family of partially entangled probes.  In the divisor
branch, a dimension converse closes the global rank-$r$ probe optimization at
$d/r$.  For rank two, the regular-polygon geometry closes the entire one-shot
list-success curve, while a simple spectral witness quantifies failure below
threshold for arbitrary full-spark ensembles and survives bounded state
perturbations.  The construction also supplies an infinite separation between
support-only feasibility tests and actual zero-error measurements.  Natural
open problems are to close the nondivisible probe optimum, obtain exact
subthreshold curves beyond rank two, and find other represented matroids whose
physical-space factor cone admits a closed positive identity resolution.

\paragraph{AI assistance disclosure.}
OpenAI Codex assisted with theorem exploration, literature triage, proof
stress testing, verifier development, and manuscript editing.  The author
selected and reviewed the claims, proofs, citations, and computational
evidence and remains responsible for the work.  AI output is not treated as
independent scientific evidence.

\begingroup
\fontsize{8.5}{10}\selectfont
\let\oldthebibliography\thebibliography
\renewcommand{\thebibliography}[1]{%
  \oldthebibliography{#1}\setlength{\itemsep}{0pt}}
\bibliographystyle{unsrt}
\bibliography{references}
\endgroup

\end{document}
